\documentclass[sigconf,review,anonymous]{acmart}

\settopmatter{printacmref=true}
\renewcommand\footnotetextcopyrightpermission[1]{}
\pagestyle{plain}

\usepackage{booktabs}
\usepackage{xspace}
\usepackage{enumitem}
\usepackage{listings}

\lstdefinestyle{paperlisting}{
  basicstyle=\ttfamily\small,
  columns=fullflexible,
  frame=single,
  breaklines=true,
  keepspaces=true,
  showstringspaces=false
}
\lstset{style=paperlisting}

\newcommand{\system}{Persistent KV Anti-Caching\xspace}
\newcommand{\kv}{KV\xspace}
\newcommand{\ttft}{TTFT\xspace}
\newcommand{\prepass}{PrePass\xspace}

\title{Persistent KV Anti-Caching for Multi-Turn Long-Context LLM Serving}

\author{Anonymous Authors}
\affiliation{%
  \institution{Anonymous Institution}
  \country{}
}

\begin{CCSXML}
<ccs2012>
 <concept>
  <concept_id>10003033.10003058.10003059</concept_id>
  <concept_desc>Networks~Network performance analysis</concept_desc>
  <concept_significance>300</concept_significance>
 </concept>
 <concept>
  <concept_id>10010520.10010521.10010537</concept_id>
  <concept_desc>Computer systems organization~Distributed architectures</concept_desc>
  <concept_significance>300</concept_significance>
 </concept>
 <concept>
  <concept_id>10010147.10010178.10010179</concept_id>
  <concept_desc>Computing methodologies~Natural language processing</concept_desc>
  <concept_significance>300</concept_significance>
 </concept>
</ccs2012>
\end{CCSXML}

\ccsdesc[300]{Networks~Network performance analysis}
\ccsdesc[300]{Computer systems organization~Distributed architectures}
\ccsdesc[300]{Computing methodologies~Natural language processing}

\keywords{KV cache, long-context serving, anti-caching, LLM serving, tiered storage}

\begin{abstract}
Long-context LLM serving is increasingly used in enterprise RAG, code
assistance, and multi-agent workflows, where long historical contexts are
revisited across many turns. Recent serving systems improve KV reuse through
in-engine prefix caching, external KV stores, and SSD-backed data movement.
However, they still leave a critical state-management problem underexplored:
when reusable historical KV exceeds GPU memory, a cold KV object is persistent
but not execution-ready, and discovering it during prefill or decode can turn
storage I/O into online tail latency.

This paper introduces \system, a memory-primary KV management architecture for
exact multi-turn long-context serving. The key idea is to treat cold historical
KV as anti-cached state rather than as a conventional cache block. \system is
developed with three designs: 1) a KV Evicted Index that records committed KV
ranges across GPU-ready, CPU-ready, SSD-cold, fetching, loading, missing, and
mismatched states; 2) a KV PrePass that enumerates required historical KV before
a request enters the LLM engine and enforces a ready-barrier admission rule; and
3) a packed cold-object layout that stores cold KV as offset-addressable
sequential objects instead of many per-layer files.

We implement a vLLM-based prototype with a sidecar control plane and a custom
KV connector. Preliminary experiments on Qwen2.5-14B-Instruct show that, when
PrePass restore is hidden before online serving, \system reduces online TTFT to
0.447x, 0.295x, and 0.161x of full prefill at 2K, 8K, and 16K reusable prefixes,
respectively. On a LongMemEval-S sample, it achieves 0.537x and 0.297x online
TTFT at 2K and 8K. These results demonstrate the feasibility of exact persistent
KV reuse without synchronous SSD misses, while also showing that
restore-inclusive SLA guarantees and production-scale 1M-token serving remain
open engineering problems.
\end{abstract}

\begin{document}

\maketitle

\section{Introduction}

Long-context LLM serving is becoming a stateful systems problem. In addition to
the quadratic or near-quadratic cost of attention kernels, production-like
workflows repeatedly prefill long text that has already been processed in
previous turns. Enterprise RAG applications may attach the same organizational
policy or retrieved documents to many requests. Code assistants may revisit the
same repository files while the user asks small follow-up questions. Multi-agent
workflows may carry shared system prompts, tool descriptions, intermediate
plans, and long-lived memory across turns. In these cases, historical context is
valuable because it is expensive to recompute, but it is also too large and too
numerous to keep entirely in GPU memory.

Modern serving systems already provide important pieces of the answer.
PagedAttention and vLLM reduce fragmentation and enable efficient in-memory KV
sharing within the serving engine~\citep{kwon2023pagedattention}. Automatic
prefix caching can skip prefill for hot prefixes in the same cache
lifetime~\citep{vllmPrefixCaching}. External KV systems such as LMCache extract
KV from the engine and enable reuse across requests and
engines~\citep{liu2025lmcache}. Disaggregated systems such as Mooncake make KV
a first-class scheduling object~\citep{qin2024mooncake}. Recent SSD-backed and
multi-tier systems, including Tutti and KVDrive, directly target the data
movement bottleneck between GPU memory, host memory, and
SSD~\citep{qiu2026tutti,lin2026kvdrive}. These systems establish that KV is now
a storage and scheduling problem, not only an attention-kernel optimization
problem.

However, multi-turn long-context serving exposes a distinct problem that is not
fully captured by either hot prefix caching or faster SSD transfer:
\emph{how should a serving system manage cold historical KV as persistent state
objects while preserving exactness and protecting online latency?} A cached
prefix block in HBM is immediately useful; a KV object on SSD or 3FS is not. It
must be found, checked against the current model and prompt state, restored into
an execution tier, and admitted before attention depends on it. If the engine
discovers this requirement during prefill or decode, the request either stalls
on storage I/O or falls back to recomputation. Conversely, blindly prefetching
all possible histories wastes bandwidth and pollutes memory. The missing
abstraction is therefore not only a faster KV data path, but a control-plane
state model that decides when cold historical KV is allowed to become execution
state again.

This paper revisits database anti-caching for persistent KV serving. In
database anti-caching, main memory is the primary execution storage, cold tuples
are evicted to disk, and an in-memory evicted table records where each tuple
lives~\citep{debrabant2013anticaching}. When a transaction touches evicted data,
the system enters a pre-pass phase to discover the needed disk blocks, fetches
them asynchronously, and restarts the transaction after the data is available.
The same principle fits historical KV reuse with one crucial change: an LLM
engine cannot speculatively execute with missing historical KV because
attention semantics would be incomplete. Therefore, the pre-pass must be a
control-plane procedure before the request enters the GPU execution path.

We propose \system, a system design that adapts anti-caching to exact KV reuse.
The design is guided by four principles.

\textbf{First, separate execution tiers from persistence tiers.} HBM and host
DRAM are the tiers from which KV can become execution-ready; SSD and 3FS are
cold capacity tiers. A KV range in SSD is not a slow ready object. It is a
non-ready object that must be restored, verified, and admitted.

\textbf{Second, keep cold-state metadata in memory.} A KV Evicted Index records
committed historical KV ranges with correctness keys, token ranges, layer
ranges, tier states, object identifiers, offsets, lengths, checksums, and
in-flight restore or load identifiers. This lets the sidecar classify required
historical KV without scanning storage or asking the GPU to discover misses.

\textbf{Third, expose cold misses before admission.} KV PrePass tokenizes an
incoming request, computes a correctness key, resolves session lineage or prefix
candidates, enumerates required historical ranges, and classifies them as
GPU-ready, CPU-ready, SSD-cold, fetching, loading, missing, or mismatched. The
admission rule is explicit: requests with required cold KV are delayed,
requeued, or fallbacked; they are not allowed to synchronously read SSD in the
execution path. This transforms a cold miss from an opaque TTFT/ITL tail event
into a schedulable restore operation.

\textbf{Fourth, use coarse physical objects with fine logical metadata.} A long
prefix for a multi-layer transformer may otherwise become dozens or hundreds of
small files. That layout is simple but fragile at 8K, 16K, 32K, and beyond,
because file count, metadata overhead, Python copy overhead, and tail latency
dominate. Packed cold objects provide a block-table-like representation: the
cold tier stores large sequential objects, while the in-memory index stores the
offset table needed to restore relevant ranges.

Our current prototype is deliberately narrow. It is a vLLM-integrated research
prototype, not a production 1M-token serving system. It implements the sidecar
control plane, correctness-keyed manifests, a vLLM KV connector, true cold-tier
file migration, a packed cold-object backend, a PrePass endpoint, a ready
barrier, and a LongMemEval-S workload adapter. The most important evidence is
semantic and directional. The prototype can persist KV across vLLM restart,
demote it to a cold object, reject online execution while the object is cold,
restore it through PrePass, verify it, and then load exact external KV during a
reuse request with zero synchronous SSD misses. In controlled 2K/8K/16K
experiments, the online portion of TTFT improves as reusable context grows. But
when restore is charged to the arriving request, the system only approaches
break-even at 16K in our local setup. This limitation clarifies the central
research agenda: \system is valuable when workflows provide lead time, idle
windows, or predictable next-state transitions that allow restore to be hidden,
and when the restore/load data path is engineered to avoid cold-tier tail
latency.

In summary, this paper makes the following contributions:
\begin{itemize}[leftmargin=*]
  \item We formulate persistent historical KV serving as an anti-caching problem
  rather than a conventional cache-hit problem.
  \item We define a multi-tier KV state model and an in-memory KV Evicted Index
  that separates GPU-ready, CPU-ready, SSD-cold, fetching, loading, missing,
  and mismatched states.
  \item We design KV PrePass and a ready-barrier admission protocol that
  prevents synchronous cold-tier misses from entering prefill or decode.
  \item We introduce packed cold objects as the cold-tier data model for exact,
  verifiable, offset-addressable KV restoration.
  \item We report preliminary vLLM prototype evidence on synthetic scaling
  prompts and LongMemEval-S, while identifying the remaining gap to a
  production-scale systems evaluation.
\end{itemize}

\section{Related Work}

\subsection{Database Anti-Caching}

\system is directly inspired by DeBrabant et al.'s anti-caching architecture
for main-memory OLTP databases~\citep{debrabant2013anticaching}. Anti-caching
reverses the traditional buffer-pool hierarchy: main memory is the primary
storage, and cold tuples are moved to disk when memory is exhausted. The system
maintains an in-memory Evicted Table that maps evicted tuples to disk blocks, a
Block Table for disk-resident blocks, and a pre-pass phase that discovers
evicted blocks needed by a transaction before fetching them together. The core
lesson for KV serving is not the particular tuple layout, but the control
invariant: cold data should be explicitly indexed, fetched outside the critical
execution path, and merged back into the execution tier before the transaction
proceeds.

However, KV serving violates two assumptions behind the original database
setting. First, the evicted object is not a tuple but a model-, tokenizer-,
position-, layout-, and prompt-specific KV state. A correctness key is therefore
part of the object identity. Second, an LLM request cannot execute a meaningful
``trial'' decode with missing historical KV. Different from database
anti-caching, \system moves the pre-pass into the serving control plane, where
required KV is resolved from session lineage and prefix manifests before vLLM
execution begins.

\subsection{In-Engine KV Memory Management and Prefix Caching}

PagedAttention and vLLM provide the foundation for efficient KV memory
management in modern LLM serving~\citep{kwon2023pagedattention}.
PagedAttention uses paging-inspired block management to reduce KV fragmentation
and support flexible sharing of KV cache within and across requests. vLLM's
automatic prefix caching further hashes computed blocks and reuses matching
prefixes in the engine cache~\citep{vllmPrefixCaching}. These mechanisms are
highly effective for hot prefixes within the engine's cache lifetime.

In contrast, \system targets historical KV after it leaves the hot in-engine
cache. Such KV may survive process restart, capacity eviction, and cold-tier
demotion, and it must be rediscovered and admitted before reuse. In our
prototype, strict restart baselines remain close to full prefill at 2K, 8K, and
16K, while B5 external reuse can still load historical KV after restart. Thus,
automatic prefix caching is the right hot-cache baseline, whereas \system
addresses the persistent cold state below it.

\subsection{External KV Reuse and Enterprise KV Cache Layers}

LMCache is the closest practical external-cache system: it extracts and stores
KV caches outside GPU memory, supports vLLM and SGLang, and exposes a control
API for cache orchestration across GPU, CPU, storage, and network
layers~\citep{liu2025lmcache}. It is especially relevant to enterprise
scenarios such as multi-round QA and document analysis, where KV reuse can
reduce TTFT and GPU cycles.

\system is complementary but more specific. We do not claim that externalizing
KV is new. Our focus is the admission semantics for cold historical KV: an
in-memory residency index, explicit GPU/CPU/SSD states, a PrePass that
enumerates required historical ranges before execution, a ready barrier, and a
policy that forbids synchronous cold-tier misses on the online path. These
mechanisms are intended to make external KV reuse predictable under SLA
constraints, not merely possible.

CacheBlend and related RAG-oriented KV reuse systems address a different
correctness challenge: precomputed chunks are often not exact prefixes, so they
must be fused or partially recomputed to preserve quality. Our current design
chooses the exact-prefix and committed-history regime. It does not attempt to
reuse arbitrary non-prefix chunks or approximate attention; instead it asks how
to persist, schedule, and restore exact historical KV
efficiently~\citep{yao2024cacheblend}.

\subsection{Disaggregated and KV-Centric Serving}

Mooncake proposes a KVCache-centric disaggregated architecture that separates
prefill and decode clusters and uses underutilized CPU, DRAM, and SSD resources
to manage KV across the serving system. It also treats SLO-aware scheduling and
early rejection as first-class mechanisms~\citep{qin2024mooncake}. This is a
strong signal that KV management has become a system-wide scheduling problem.

\system can be viewed as a cold-state subsystem that could live inside such a
KV-centric architecture. Its distinctive mechanism is the anti-cache transition:
committed historical KV moves from execution tiers to a persistent cold tier,
remains discoverable through an in-memory index, and returns to execution tiers
through PrePass and admission. We do not replace disaggregated scheduling.
Instead, we specify how cold historical state should be represented and admitted
when it is no longer hot in the serving engine.

\subsection{SSD-Backed and Multi-Tier KV Systems}

Tutti and KVDrive are especially close to our target. Tutti focuses on making
SSD-backed KV practical by reworking the SSD-HBM data path around GPU-centric
object abstractions, GPU direct object I/O, and slack-aware
scheduling~\citep{qiu2026tutti}. KVDrive builds a holistic multi-tier KV cache
management system spanning GPU memory, host DRAM, and SSD, with
attention-aware placement, pipeline scheduling, and cross-tier coordination for
long-context inference~\citep{lin2026kvdrive}. These systems directly address
the fact that KV movement can dominate long-context serving latency.

Our contribution should not be stated as ``we add SSD to KV cache'' or ``we
manage GPU/DRAM/SSD tiers''; KVDrive already makes that framing too broad. The
narrower claim is that multi-turn serving needs a persistent anti-cache
semantics for historical KV. KVDrive primarily motivates how to move critical
KV efficiently across tiers during long-context inference. \system asks when a
cold historical KV object is allowed to become execution state again: it must
be exact, correctness-key matched, indexed, restored before admission, and
observed by a ready barrier. In this sense, Tutti and KVDrive are strong
data-plane and multi-tier baselines, while our design emphasizes the
control-plane semantics for persistent historical reuse.

\subsection{Long-Term Memory Workloads}

LongMemEval evaluates long-term interactive memory in chat assistants, with
questions spanning information extraction, multi-session reasoning, temporal
reasoning, knowledge updates, and abstention~\citep{wu2024longmemeval}. Its
multi-session histories map well to our workload abstraction:
\texttt{haystack\_sessions} become committed historical KV prefixes, and the
current question becomes the online suffix. We use LongMemEval-S as an early
public workload because it captures repeated, long-lived conversational state
more realistically than synthetic repeated tokens. Our current use is
systems-level rather than accuracy-level: we measure whether historical prefixes
can be stored, demoted, restored, and reused with real multi-turn text.
Answer-quality evaluation and larger workload coverage remain future work.

\section{Design}

\subsection{Design Goals}

The design of \system is guided by six requirements that separate persistent
state management from conventional cache-hit optimization.

\paragraph{Exact reuse.} A restored KV range must be equivalent to the KV that
the same model would have produced by full prefill over the same token sequence
under the same model revision, tokenizer, position encoding, dtype, KV layout,
attention backend, and adapter configuration. The system does not use lossy KV
compression or approximate sparse attention in its first design.

\paragraph{Execution-tier clarity.} GPU HBM and host DRAM are execution tiers.
SSD and 3FS are cold persistence tiers. A cold-tier object may be durable and
reusable, but it is not ready for execution until restored and verified.

\paragraph{No synchronous cold miss.} Online prefill and decode must not block
on unexpected SSD/3FS reads for required historical KV. Cold misses must appear
as admission decisions, restore tasks, or fallbacks.

\paragraph{Multi-turn persistence.} Historical KV should survive the lifetime
of a single request and, in the target design, the lifetime of a single engine
process. This supports restarted engines, capacity pressure, cross-engine
serving, and idle session restoration.

\paragraph{Schedulability.} Required cold bytes, expected restore time, H2D load
time, and queue state should be visible before admission. The system should be
able to delay, requeue, fallback, or reject low-locality requests rather than
hiding them inside tail latency.

\paragraph{Workload boundary honesty.} The target workload is high-reuse
historical context, not arbitrary low-locality 1M-token prompts. Low-reuse
requests should be measured as negative controls and handled through separate
queues or fallback paths.

\subsection{Non-Goals}

The first design does not modify attention kernels, change model semantics, or
claim that random 1M-token requests can be served within tight TTFT by storage
alone. It also does not claim that the current prototype already reaches 1M
tokens or production 3FS performance. The intended contribution is a systems
architecture and prototype evidence for exact, persistent, and schedulable KV
reuse.

\subsection{System Overview}

\system sits between the application gateway and the LLM engine, as a
control-plane layer for historical KV state. The sidecar owns durable metadata,
including session lineage, prefix manifests, correctness keys, the KV Evicted
Index, cold-object manifests, restore queues, and metrics. The vLLM connector
owns controlled data movement into and out of the engine's KV representation.

The request path has five stages:
\begin{enumerate}[leftmargin=*]
  \item \textbf{Commit.} After a request computes a reusable prefix, the
  connector stores the KV tensors and the sidecar commits a manifest with token
  ranges, layer ranges, block alignment, correctness key, checksums, and
  version metadata.
  \item \textbf{Demotion.} When hot memory must be reclaimed, the system packs
  selected historical ranges into cold objects on SSD or 3FS, verifies
  checksums, and updates the index from ready execution state to
  \texttt{SSD\_COLD}.
  \item \textbf{PrePass.} Before a future request enters vLLM, the sidecar
  enumerates its required historical ranges and classifies each range with the
  in-memory index.
  \item \textbf{Restore and admission.} If all required ranges are
  execution-ready, the request is admitted. If some ranges are cold, restore
  tasks are issued and the request is delayed or requeued. Admission occurs
  only after the ready barrier observes that required ranges are restored and
  correctness-checked.
  \item \textbf{External load and suffix prefill.} The connector loads
  historical KV and vLLM computes only the new suffix and decode tokens.
\end{enumerate}

This architecture intentionally separates policy from engine internals. The
sidecar determines whether a request is allowed to use external KV, while the
connector performs the actual load/store under vLLM's scheduling lifecycle.

\subsection{KV Object Model and Correctness Key}

A KV object is a committed, versioned range of transformer KV states. The
logical unit is:
\begin{lstlisting}
KVRange = {
  tenant_id,
  session_id,
  prefix_id,
  token_start,
  token_end,
  layer_start,
  layer_end,
  block_ids,
  correctness_key,
  committed_epoch
}
\end{lstlisting}

The correctness key is part of the object identity. It includes model identity
and revision, tokenizer hash, prompt token hash, position encoding and scaling
configuration, dtype, KV layout, attention backend, block size, adapter state,
and tenant or security domain. A prefix identifier alone is not sufficient: two
objects with the same text but different tokenizer, RoPE scaling, LoRA adapter,
or KV layout are different objects and must not be reused.

The first prototype stores whole-prefix, all-layer ranges because that is
sufficient to validate the control path. The design leaves room for smaller
token extents and layer groups, which become important for partial promotion
and parallel restore.

\subsection{KV Evicted Index}

The KV Evicted Index is the in-memory metadata structure that makes
anti-caching possible. It plays the same role as an evicted table in a
database: the data itself may be cold, but the system must always know where it
is and whether it can be used.

Each entry records:
\begin{lstlisting}
KVIndexEntry = {
  prefix_id,
  token_start,
  token_end,
  correctness_key_hash,
  tier,
  ready,
  object_id,
  cold_uri,
  extents[],
  checksum,
  size_bytes,
  last_access_epoch,
  access_count,
  fetching_request_id,
  loading_request_id
}
\end{lstlisting}

The index classifies ranges into seven states, as shown in
Table~\ref{tab:states}.

\begin{table*}[t]
  \centering
  \caption{Multi-tier KV states used by the KV Evicted Index.}
  \label{tab:states}
  \begin{tabular}{lll}
    \toprule
    State & Meaning & Admission Implication \\
    \midrule
    \texttt{GPU\_READY} & KV is already in HBM & Can enter execution directly \\
    \texttt{CPU\_READY} & KV is in host DRAM and verified & Needs CPU-to-GPU load budget or connector load \\
    \texttt{SSD\_COLD} & KV is only in SSD/3FS cold tier & Must restore before admission \\
    \texttt{FETCHING} & SSD/3FS-to-CPU restore is in progress & Do not enqueue duplicate restore \\
    \texttt{LOADING} & CPU-to-GPU load is in progress & Wait for load completion or load barrier \\
    \texttt{MISSING} & No reusable KV range exists & Full prefill, fallback, or reject \\
    \texttt{MISMATCH} & Correctness key does not match & Reuse is forbidden \\
    \bottomrule
  \end{tabular}
\end{table*}

The split between \texttt{CPU\_READY} and \texttt{GPU\_READY} is essential.
Treating DRAM-ready KV as execution-ready would hide H2D cost inside the online
path and understate TTFT. In the current prototype, PrePass already reports
load-readiness fields such as CPU-ready token count, expected H2D bytes, and
estimated H2D load time; a production system should close the loop by writing
connector load-completion events back into the index.

\subsection{KV PrePass}

KV PrePass is the anti-cache probe for LLM serving. It runs before a request is
submitted to the engine. Its inputs are the request tokens, session or prefix
candidates, and the desired correctness key. Its output is a reuse plan, a
restore plan, and an admission decision.

PrePass performs:
\begin{enumerate}[leftmargin=*]
  \item Tokenize the prompt and compute the correctness key.
  \item Resolve session lineage or shared-prefix candidates.
  \item Enumerate required historical ranges and the new suffix range.
  \item Look up required ranges in the KV Evicted Index.
  \item Return classification sets and byte estimates.
  \item Enqueue restore for \texttt{SSD\_COLD} ranges, avoid duplicate enqueue
  for \texttt{FETCHING}, and block reuse for \texttt{MISSING} or
  \texttt{MISMATCH}.
  \item Emit \texttt{ADMIT}, \texttt{DELAY},
  \texttt{FULL\_PREFILL\_FALLBACK}, or \texttt{REJECT}.
\end{enumerate}

The key policy is simple: a request with required cold KV cannot enter the GPU
fast path. It may be delayed until restore completes; it may execute full
prefill if recomputation is cheaper or the SLA cannot tolerate waiting; or it
may be rejected from the strict queue. Thus, PrePass turns cold-tier uncertainty
into an explicit scheduling decision instead of a hidden TTFT or decode-tail
spike.

\subsection{Admission and Ready Barrier}

The ready barrier is the enforcement point for PrePass. It checks that every
required range in the reuse plan is present, verified, and loadable before
decode depends on it. In the prototype, the barrier is implemented through
sidecar admission state and connector metadata. A cold probe first returns
\texttt{DELAY} with \texttt{required\_kv\_not\_ready\_before\_decode}; after
\texttt{/prefetch/advance} completes restore and checksum verification, a
repeated admission returns \texttt{ADMIT}.

This barrier also provides a clean metric surface:
\begin{itemize}[leftmargin=*]
  \item \texttt{sync\_ssd\_miss\_total} should remain zero for admitted
  strict-SLA requests;
  \item \texttt{ready\_barrier\_all\_ready} should be true before external
  load;
  \item \texttt{external\_load\_observed} confirms the engine used the
  connector path;
  \item restore-inclusive TTFT and online TTFT must be reported separately.
\end{itemize}

The distinction between online TTFT and restore-inclusive TTFT is central. If a
workflow can predict that a historical object will be needed, restore can
happen before the user-visible request. If restore starts only when the request
arrives, the restore cost must be charged to the request. The system should not
hide this distinction.

\subsection{Packed Cold Objects}

Early prototypes often store one KV tensor file per layer. That layout is
manageable for tiny prompts, but it becomes a poor cold-tier format as context
grows. A 32-layer or 48-layer model with many prefixes quickly creates many
small files, fragmented reads, repeated metadata operations, and long restore
tails. The anti-cache analogue is the difference between tuple identifiers and
disk blocks: the system needs fine-grained logical metadata but coarse-grained,
sequential physical I/O.

A packed cold object stores many layer files or extents in one object:
\begin{lstlisting}
packed_object.bin
packed_manifest.json
\end{lstlisting}

The manifest records the object id, correctness key hash, token range, layer
range, layout version, checksum, and an offset table:
\begin{lstlisting}
KVExtentRef = {
  layer_name,
  token_start,
  token_end,
  offset,
  length,
  checksum
}
\end{lstlisting}

The current \texttt{packed\_v1} prototype uses a layer-major packed object
because it was the lowest-risk path to validate true cold restore. Even this
first version surfaced an important systems lesson: packing is not
automatically faster. The initial implementation reduced file count but was
slower because it used whole-object reads, byte slicing, and repeated checksum
passes. After changing demotion and restore to streaming copy paths and
avoiding redundant whole-object summaries, packed restore became slightly
faster than the local per-layer path in a 2K/8K tiny-profile benchmark. This
reinforces the design direction while also showing that a full systems
evaluation must include layout and data-path engineering, not just metadata
design.

The next design step is extent-major packing, where token extents become the
outer layout and layers are stored within each extent. Such a layout would let
PrePass restore only the ranges required by a request and align better with
partial promotion.

\subsection{Restore Scheduler and Multi-Tier Transitions}

\system uses explicit state transitions:
\begin{lstlisting}
GPU_READY -> CPU_READY -> SSD_COLD
SSD_COLD -> FETCHING -> CPU_READY
CPU_READY -> LOADING -> GPU_READY
\end{lstlisting}

Demotion moves ranges out of execution tiers when they become cold. Restore
moves them back from SSD/3FS to CPU. Load moves them from CPU to GPU or makes
them available to the connector during execution. The scheduler must account for
at least four resources: cold-tier bandwidth, host memory capacity, H2D
bandwidth, and GPU scheduling slack.

The current prototype provides a deterministic queue and an asynchronous queue
interface. It estimates bytes and restore time, prevents duplicate restore for
already-fetching ranges, exposes queue metrics, and records actual restore bytes
and checksums. It does not yet implement a production-grade executor with queue
depth tuning, pinned staging buffers, H2D overlap, or 3FS-specific semantics.
These are engineering gaps in the data path, not changes to the anti-caching
state model.

\subsection{vLLM Integration}

The prototype integrates with vLLM through a sidecar and a custom connector.
The sidecar translates admission results into \texttt{kv\_transfer\_params};
the connector receives metadata during scheduling and worker execution, stores
KV for committed prefixes, and loads external KV for reuse requests. We use
restart experiments to distinguish true external persistence from vLLM's
in-process prefix cache.

The connector path revealed two practical issues. First, reuse requests must be
load-only unless explicitly asked to commit a new prefix. An earlier version
accidentally wrote back full KV during reuse, inflating TTFT by seconds. Adding
a \texttt{store\_policy} field fixed this issue and reduced 2K B3/B5 TTFT from
multi-second artifacts to hundreds of milliseconds. Second, long-context store
experiments must verify saved token counts. An 8K run initially saved only 2K
tokens because the engine's batched-token configuration exposed an incomplete
prefill to the connector. The prototype now rejects rows where saved tokens do
not match the requested prefix.

These lessons are part of the system design. Persistent KV reuse needs
correctness guards around connector lifecycle, block alignment, token ranges,
and store/load policies; otherwise, an apparent reuse path may be measuring
connector artifacts rather than KV state reuse.

\subsection{Prototype Evaluation Scope}

The prototype currently supports two evaluation modes. The synthetic mode uses
controlled prefix lengths to expose scaling. The LongMemEval-S mode uses a
public multi-session history to replace repeated synthetic tokens with real
conversation text.

On synthetic 2K/8K/16K prefixes with Qwen2.5-14B-Instruct, B5 online TTFT is
0.447x, 0.295x, and 0.161x of full prefill when restore is completed before
the online request. The same runs show that restore-inclusive latency is still
1.708x and 1.576x of full prefill at 2K and 8K, and approximately breaks even
at 16K. On LongMemEval-S, one real sample shows online/B0 ratios of 0.537x at
2K and 0.297x at 8K, with restore-inclusive ratios of 1.636x and 1.544x.

These results support two claims and reject one premature claim. They support
the claim that exact persistent KV can be restored before admission and loaded
without synchronous SSD misses. They also support the directional claim that
online TTFT savings grow with reusable history length. They do not support a
claim that the current prototype already satisfies a 1M-token production SLA or
outperforms strong baselines under all accounting methods.

\section{Discussion and Open Problems}

The strongest version of this paper will require a broader evaluation. The
system must compare against vLLM APC hot and restart cases, LMCache, naive
synchronous cold restore, Tutti-style SSD data paths, Mooncake-style
disaggregated scheduling, and KVDrive-style multi-tier management. It must also
report p95/p99 latency, concurrency, bandwidth utilization, admitted-workload
historical byte hit rate, and negative-control workloads. The current prototype
should therefore be understood as establishing the design invariant and exposing
the next performance bottlenecks: cold restore executor time and
CPU-ready-to-GPU load time.

The core research bet remains promising. KVDrive and Tutti show that KV data
movement is a first-order bottleneck; LMCache and Mooncake show that KV must be
managed outside a single GPU cache; LongMemEval-style workloads show that
long-lived histories are real. \system adds a state-management semantics to this
design space: cold historical KV should be persistent, exact, indexed, restored
before admission, and protected by a ready barrier. That is the system problem
this project should continue to sharpen.

\section*{AI Disclosure}

This paper was prepared with the assistance of AI-powered academic writing
tools. The content, arguments, and conclusions were directed and reviewed by
the author(s), who take responsibility for the accuracy and integrity of this
work.

\bibliographystyle{ACM-Reference-Format}
\bibliography{references}

\end{document}
